\chapter*{第六部分收束：从项目交付到课程资产}
\addcontentsline{toc}{chapter}{第六部分收束：从项目交付到课程资产}

Part VI 的章节很多，原因不是 capstone 项目本身应该变得复杂，而是一个可靠项目从来不只是一份最终报告。它同时面对学生、助教、教师、下一轮课程团队和外部采用者。不同读者关心的问题不同：学生关心怎样完成项目，助教关心怎样公平评分，教师关心课程能否运行，维护者关心下一学期能否接手，外部合作者关心边界是否清楚。

因此，读完 Part VI 后，最重要的收获不是记住每一张表叫什么，而是理解一条主线：\textbf{capstone 的目标，是把一次学生分析变成别人可以复查、评价、交接和维护的课程资产。}

\section*{三层交付模型}

可以把 Part VI 压缩成三层交付模型。

第一层是\textbf{学生项目包}。它回答“这个项目本身是否可信”：

\begin{itemize}
    \item 科学问题是否具体；
    \item 数据来源、许可证、字段、单位和处理步骤是否说清；
    \item baseline、模型、评估和错误分析是否完整；
    \item notebook、报告、图表和展示是否能互相支撑；
    \item AI 使用记录和人工签字边界是否明确。
\end{itemize}

第二层是\textbf{课程运行包}。它回答“一门课能否稳定带学生完成这些项目”：

\begin{itemize}
    \item 选题、proposal、pilot、midpoint review 和 final delivery 是否有清楚节奏；
    \item rubric、TA guide、学生 handout 和试教反馈是否互相一致；
    \item 课程日历是否提前暴露风险，而不是把风险留到期末；
    \item 反馈、修订、归档和公开边界是否能形成闭环。
\end{itemize}

第三层是\textbf{长期维护包}。它回答“这套课程资源能否被下一轮、下一位教师或外部采用者继续使用”：

\begin{itemize}
    \item instructor handoff 是否让新教师知道从哪里开始；
    \item public release index、adoption note 和 maintenance cadence 是否说明支持边界；
    \item reboot preflight、failure-mode playbook 和 substitute handoff 是否覆盖运行中断；
    \item shutdown warm-start、mentor relay、community memory 和 guest intake 是否把跨届经验保存下来。
\end{itemize}

\section*{Part VI 的共同算法}

尽管每一章的对象不同，它们背后其实反复使用同一条算法：

\begin{enumerate}
    \item 先定义当前材料要服务谁；
    \item 再检查能否复现、评分、发布或交接；
    \item 优先处理边界、责任人和证据缺口；
    \item 把材料 route 到 ready、revise、review、assign owner 或 block；
    \item 对非 ready 项写下下一步动作，而不是只给一个模糊评价。
\end{enumerate}

这条算法故意很朴素，因为课程运行中真正危险的常常不是复杂问题，而是简单问题没人负责：链接坏了、owner 缺失、许可不清、rubric 没对齐、notebook 没跑通、学生不知道 AI 使用边界。Part VI 的作用，就是把这些容易被忽视的地方提前变成可检查的对象。

\section*{从学生视角看，哪些必须完成}

如果你是学生，第一次完成 capstone 时不需要深入掌握所有课程维护章节。但至少要交出下面五类证据：

\begin{itemize}
    \item 一个能解释清楚的问题和范围；
    \item 一份能复现的数据与 notebook 入口；
    \item 一个 baseline 和一个经过诊断的主要模型；
    \item 一组结果图、错误案例和限制说明；
    \item 一份报告、展示、引用记录和 AI-use statement。
\end{itemize}

这些材料对应的是第 39--42 章的核心要求。后面的课程运行与维护章节，更多是在解释为什么教师和助教会要求你留下这些证据：不是为了增加负担，而是为了让你的项目真正能被别人读懂、复查和评价。

\section*{从教师视角看，哪些最容易失效}

如果你是教师或课程维护者，Part VI 最值得反复检查的是四类失效点：

\begin{itemize}
    \item \textbf{入口失效}：学生或外部读者找不到正确 notebook、数据、rubric 或提交入口；
    \item \textbf{责任失效}：每个问题都重要，但没有明确 owner；
    \item \textbf{边界失效}：数据、版权、隐私、AI 使用或外部合作边界没有写清；
    \item \textbf{记忆失效}：上一轮课程学到的经验没有进入 handoff、ledger 或 warm-start。
\end{itemize}

这些失效通常不会在第一天显现，却会在期末、换教师、公开发布或外部采用时集中爆发。Part VI 后半部分看似“运维”，其实是在保护课程的长期可用性。

\section*{进入全书收束前}

到这里，全书的技术线已经闭合：从 Linux、Git、数据、图表，到机器学习、深度学习、现代 AI，再到 capstone 交付。最后需要带走的不是更多章节，而是一种工作方式：任何科学结果都应该能回答“数据从哪里来、代码怎样跑、模型怎样评估、错误怎样解释、别人怎样复查、下一轮怎样接手”。

如果这六个问题能回答，一个本科 capstone 就已经不只是作业，而是一个小型、可信、可维护的科研训练成果。
